\chapter{Introducción}

La locomoción de robots cuadrúpedos ha sido ampliamente estudiada debido a su capacidad para desplazarse sobre terrenos irregulares y mantener la estabilidad ante perturbaciones. Plataformas como MIT Cheetah~3 y ANYmal muestran que una arquitectura mecánica articulada, combinada con estimación de estado y control de cuerpo completo, permite ejecutar desplazamientos robustos en condiciones experimentales exigentes~\cite{bledt2018cheetah3,hutter2016anymal}.

El desempeño locomotor depende de la configuración estructural, la posición del cuerpo, las trayectorias de los pies y las fuerzas de contacto. Por esta razón, el modelado cinemático y dinámico, el cálculo de las referencias articulares y la coordinación de las fases de apoyo y oscilación son elementos centrales del diseño de la marcha~\cite{zamani2011dynamics,bellicoso2018dynamic}.

Los enfoques modernos incorporan el estado del robot y la interacción pie--suelo para ajustar el movimiento. El control predictivo basado en modelo y la optimización de cuerpo completo permiten imponer restricciones de contacto y de los actuadores, además de responder a perturbaciones durante la locomoción~\cite{dicarlo2018mpc,neunert2018wholebody}. Estos resultados evidencian que el desarrollo de un sistema cuadrúpedo requiere integrar de manera coherente el diseño mecánico, el modelo matemático, la simulación y la arquitectura electrónica de control.

El presente anteproyecto propone una ruta integral para desarrollar el sistema de locomoción de una plataforma cuadrúpeda existente. El proceso parte de la caracterización física del robot, continúa con el modelado cinemático y dinámico, y culmina con una arquitectura de control orientada a implementar patrones básicos de locomoción, principalmente postura, paso y gateo.
